\documentclass[12pt, a4paper]{article}

% Γλώσσα και Κωδικοποίηση
\usepackage[utf8]{inputenc}
\usepackage[greek,english]{babel}
\selectlanguage{greek}
\usepackage{alphabeta}
\usepackage{textgreek}  % Για ελληνικά γράμματα σε μαθηματική λειτουργία
\newcommand{\lt}[1]{\textlatin{#1}}
\usepackage{enumitem}

% Μαθηματικά
\usepackage{amsmath}
\usepackage{amssymb}

% Μορφοποίηση Σελίδας και Διάστιχο
\usepackage{geometry}
\geometry{left=3cm, right=2cm, top=4cm, bottom=3cm}
\usepackage{setspace}
\setstretch{1.5}

% Χρώματα και Ορισμοί
\usepackage{xcolor}
\definecolor{maroon}{HTML}{AF3235}
\newcommand{\HRule}{\rule{\linewidth}{0.5mm}}

% Πακέτο για monospace κώδικα
\usepackage{verbatim}

% Εικόνες
\usepackage{graphicx}
\graphicspath{{images/}}

% Εντολή για άδεια σελίδα
\usepackage{afterpage}
\newcommand\blankpage{%
	\null
	\thispagestyle{empty}%
	\addtocounter{page}{-1}%
	\newpage}

\begin{document}
	
\begin{titlepage}
	\begin{center}
		\HRule \\[0.6cm] 
		{\huge \bfseries \strut Numerical Analysis - Lab Project 2} \\[0.4cm] 
		\HRule \\[1.5cm]
		{\large \textbf{Topic:}} \\
		{\large Numerical Integration Utility using Composite Trapezoidal and Simpson's 1/3 Rules} \\[2cm]
		
		\vfill 
		\begin{doublespacing}
			{\huge Antonios Zacharias\\}
			{\large \textbf{BSc} in Mathematics, AUTh\\}
			{\large \textbf{MSc} in Communication Networks \& System Security \\
			Computer Science Department, AUTh}
		\end{doublespacing}
	
		\vspace*{1cm} 
		\end{center}
\end{titlepage}	
	

		
% --- Part 1: Composite Trapezoidal Rule ---
\section*{1. Composite Trapezoidal Rule (trapezoid)}
\textbf{Implementation Logic}: The function \texttt{trapezoid} calculates the definite integral of a function $f$ over the interval $[a, b]$ by applying the composite trapezoidal rule.

\begin{itemize}
	\item \textbf{Algorithm}: The interval is partitioned into $n$ equal subintervals, each with a width of $h = (b-a)/n$.
	\item \textbf{Mechanism}: A \texttt{for} loop calculates the sum of the function values at all internal points from $x_1$ to $x_{n-1}$.
	\item \textbf{Calculation}: The final integral value is produced by the standard formula:
	\begin{equation*}
		I = \frac{h}{2} \left[ f(a) + f(b) + 2 \sum_{i=1}^{n-1} f(x_i) \right]
	\end{equation*}
\end{itemize}

\textbf{Execution Example}:
\begin{verbatim}
	>> f = @(x) x.*exp(2*x);
	>> trapezoid(f, 0, 4, 8)
	ans = 5764.76
\end{verbatim}

% --- Part 2: Composite Simpson's 1/3 Rule ---
\section*{2. Composite Simpson's 1/3 Rule (simpson)}
\textbf{Implementation Logic}: The function \texttt{simpson} handles numerical integration using the composite Simpson's 1/3 rule for increased precision.

\begin{itemize}
	\item \textbf{Constraint}: An \texttt{if} statement checks if $n$ is an even number; if not, the function displays an error message and terminates.
	\item \textbf{Algorithm}: The function evaluates $f(x)$ at each grid point and applies weights based on the index: 4 for odd and 2 for even.
	\item \textbf{Mechanism}: A \texttt{for} loop iterates through the internal points, using \texttt{mod(i, 2)} to identify the correct weight for each term.
	\item \textbf{Calculation}: The sum is combined with the endpoints using the formula:
	\begin{equation*}
		J = \frac{h}{3} \left[ f(a) + f(b) + \sum (\text{weighted internal points}) \right]
	\end{equation*}
\end{itemize}

\textbf{Execution Example}:
\begin{verbatim}
	>> f = @(x) x.*exp(2*x);
	>> simpson(f, 0, 4, 4)
	ans = 5670.975
\end{verbatim}

% --- Part 3: Double Integration (doubleint) ---
\section*{3. Double Integration (doubleint)}
\textbf{Implementation Logic}: The function \texttt{doubleint} acts as a modular tool to compute double integrals over rectangular domains.

\begin{itemize}
	\item \textbf{Modular Execution}: The integration is performed layer-by-layer. It defines a temporary function $g(x) = f(x, y_i)$ for each $y$ step.
	\item \textbf{Integration Flow}: 
	\begin{enumerate}
		\item It calls the \texttt{trapezoid} and \texttt{simpson} functions to solve the inner integral with respect to $x$.
		\item It then accumulates these results to perform the outer integration with respect to $y$.
	\end{enumerate}
\end{itemize}

\textbf{Execution Example}:
\begin{verbatim}
	>> f = @(x,y) x + y.^3;
	>> doubleint(f, 0, 4, 1, 3, 4, 2)
	trapezoid = 104
	simpson = 96 
	
	Examples Verified via MATLAB Command Window
\end{verbatim}



\end{document}




